\documentclass[11pt,a4paper]{article}

% =========================
% Packages
% =========================
\usepackage[utf8]{inputenc}
\usepackage[T1]{fontenc}
\usepackage[french]{babel}
\usepackage[a4paper,margin=2.2cm]{geometry}
\usepackage{lmodern}
\usepackage{amsmath,amssymb,amsthm,mathtools}
\usepackage{enumitem}
\usepackage{xcolor}
\usepackage{fancyhdr}
\usepackage{lastpage}
\usepackage{hyperref}
\usepackage{tikz}
\usetikzlibrary{arrows.meta,calc}
\usepackage{array}

% =========================
% Hyperref
% =========================
\hypersetup{
    colorlinks=true,
    linkcolor=black,
    urlcolor=blue,
    pdftitle={DM de géométrie},
    pdfauthor={},
    pdfsubject={Géométrie analytique dans le plan et dans l'espace}
}

% =========================
% General layout
% =========================
\setlength{\parindent}{0pt}
\setlength{\parskip}{0.5em}

\pagestyle{fancy}
\fancyhf{}
\fancyhead[L]{\textsc{Devoir Maison}}
\fancyhead[C]{\textsc{Géométrie analytique}}
\fancyhead[R]{\textsc{Lycée }}
\fancyfoot[C]{\thepage/\pageref{LastPage}}

% =========================
% Custom commands
% =========================
\newcommand{\R}{\mathbb{R}}
\newcommand{\vv}[1]{\overrightarrow{#1}}
\newcommand{\Boule}{\mathcal{B}}
\newcommand{\Sphere}{\mathcal{S}}
\newcommand{\Cube}{\mathcal{C}}
\newcommand{\pts}[1]{\hfill{\small\textit{(#1 pt)}}}
\newcommand{\barq}[1]{\hfill{\small\textbf{[#1 pt]}}}

\newenvironment{consigne}
{\begin{center}
\begin{minipage}{0.94\textwidth}
\small
\color{black}
\hrule
\vspace{0.4em}
\textbf{Consignes.}
}
{
\vspace{0.4em}
\hrule
\end{minipage}
\end{center}
}

% =========================
% Document
% =========================
\begin{document}

\begin{center}
    {\Large \textbf{Devoir Maison --- Géométrie analytique}}\\[0.4em]
    {\large \textbf{Raisonnement dans le plan et dans l'espace}}\\[0.8em]
    \rule{0.75\textwidth}{0.6pt}
\end{center}

\vspace{0.3em}


\vspace{0.8em}

\begin{consigne}
La rédaction doit être \textbf{soignée, rigoureuse et complète}.\\
Toute affirmation doit être \textbf{justifiée}. Un résultat de calcul seul ne suffit pas.\\
Lorsqu'une figure semble utile, on pourra en faire une, mais \textbf{aucune figure ne remplace une démonstration}.PAS D'IA, ni aucune aide extérieurs.\\
On attend une attention particulière portée aux \textbf{interprétations géométriques}.\\
\textbf{Barème indicatif :} le devoir est noté sur \textbf{40 points}.
\end{consigne}

\vspace{0.8em}

\section*{Problème 1 --- Boule dans l'espace, intersections et inclusion \hfill / 24 points}

On se place dans un repère orthonormé de l'espace.

On considère le point
\[
\Omega(1,1,1)
\]
et la boule fermée
\[
\Boule=\left\{M(x,y,z)\in\R^3 \,\middle|\, \Omega M\le 3\right\}.
\]
Sa frontière est la sphère
\[
\Sphere=\left\{M(x,y,z)\in\R^3 \,\middle|\, \Omega M=3\right\}.
\]

On considère également le cube
\[
\Cube=[0,2]\times[0,2]\times[0,2].
\]

\subsection*{Partie A --- Équation, appartenance, inclusion}

\begin{enumerate}[label=\textbf{\arabic*.}]
    \item Donner une équation cartésienne de la sphère \(\Sphere\). \barq{1}

    \item Les points
    \[
    A(4,1,1),\qquad B(1,4,1),\qquad C(1,1,-2)
    \]
    appartiennent-ils à \(\Sphere\) ? \barq{1.5}

    \item Le point
    \[
    D(2,2,2)
    \]
    appartient-il à la boule \(\Boule\) ? \barq{1}

    \item Montrer que tout point du cube \(\Cube\) appartient à la boule \(\Boule\).  
    On établira ainsi l'inclusion
    \[
    \Cube\subset \Boule.
    \]
    \barq{2}

    \item Cette inclusion est-elle stricte ? \barq{0.5}
\end{enumerate}

\subsection*{Partie B --- Intersection sphère--droite}

On considère la droite \(\Delta\) passant par \(\Omega\) et de vecteur directeur
\[
\vec u=(1,2,2).
\]

\begin{enumerate}[label=\textbf{\arabic*.},resume]
    \item Donner une représentation paramétrique de \(\Delta\). \barq{1}

    \item Déterminer l'intersection
    \[
    \Delta\cap \Sphere.
    \]
    \barq{2}

    \item Interpréter géométriquement le résultat obtenu. \barq{0.5}
\end{enumerate}

On considère maintenant la droite \(\Delta'\) définie par
\[
\begin{cases}
x=4+t\\
y=4\\
z=1
\end{cases}
\qquad (t\in\R).
\]

\begin{enumerate}[label=\textbf{\arabic*.},resume]
    \item Déterminer l'intersection
    \[
    \Delta'\cap \Sphere.
    \]
    \barq{2}

    \item Donner, dans le cas général, les différentes situations possibles pour l'intersection d'une droite et d'une sphère. \barq{1}
\end{enumerate}

\subsection*{Partie C --- Intersection sphère--plan}

On considère le plan
\[
P:\ x+y+z=3.
\]

\begin{enumerate}[label=\textbf{\arabic*.},resume]
    \item Vérifier que \(\Omega\in P\). \barq{0.5}

    \item En déduire la nature de l'intersection
    \[
    P\cap \Sphere.
    \]
    \barq{1}

    \item Préciser le centre et le rayon de cette intersection. \barq{1}
\end{enumerate}

On considère maintenant le plan
\[
Q:\ x+y+z=7.
\]

\begin{enumerate}[label=\textbf{\arabic*.},resume]
    \item Calculer la distance du point \(\Omega\) au plan \(Q\). \barq{1.5}

    \item En déduire la nature de l'intersection
    \[
    Q\cap \Sphere.
    \]
    \barq{1}
\end{enumerate}

On considère enfin le plan
\[
R:\ x+y+z=10.
\]

\begin{enumerate}[label=\textbf{\arabic*.},resume]
    \item Déterminer la nature de l'intersection
    \[
    R\cap \Sphere.
    \]
    \barq{1}

    \item Résumer les différents cas possibles pour l'intersection d'un plan et d'une sphère selon la distance du centre au plan. \barq{1}
\end{enumerate}

\subsection*{Partie D --- Intersection plan--cube}

On considère le plan
\[
H:\ x+y+z=2.
\]

\begin{enumerate}[label=\textbf{\arabic*.},resume]
    \item Montrer que ce plan coupe le cube \(\Cube\). \barq{1}

    \item Déterminer les sommets du cube appartenant au plan \(H\). \barq{1}

    \item Déterminer les points d'intersection du plan \(H\) avec les arêtes du cube. \barq{2}

    \item En déduire la nature géométrique de la section
    \[
    H\cap \Cube.
    \]
    \barq{1}

    \item Calculer l'aire de cette section. \barq{1}
\end{enumerate}

\subsection*{Partie E --- Intersection sphère--cube}

Dans toute cette partie seulement, on considère la boule fermée
\[
\Boule'=\left\{M(x,y,z)\in\R^3 \,\middle|\, \Omega M\le \sqrt3\right\}
\]
et sa frontière, la sphère
\[
\Sphere'=\left\{M(x,y,z)\in\R^3 \,\middle|\, \Omega M=\sqrt3\right\},
\]
où \(\Omega(1,1,1)\).

\begin{enumerate}[label=\textbf{\arabic*.},resume]
    \item Montrer que
    \[
    \Cube\cap \Sphere'\neq \varnothing
    \qquad\text{et}\qquad
    \Cube\subset \Boule'.
    \]
    \barq{1.5}

    \item Déterminer les sommets du cube appartenant à la sphère \(\Sphere'\). \barq{1}

    \item Déterminer les arêtes du cube qui rencontrent la sphère \(\Sphere'\). \barq{1.5}

    \item Pour chacune de ces arêtes, calculer explicitement le ou les points d'intersection avec la sphère \(\Sphere'\). \barq{2}

    \item Décrire aussi précisément que possible l'ensemble
    \[
    \Cube\cap \Sphere'.
    \]
    \barq{1}
\end{enumerate}

\vspace{1em}
\hrule
\vspace{1em}

\section*{Problème 2 --- Lieu géométrique dans le plan et structure cachée \hfill / 16 points}

On se place maintenant dans un repère orthonormé du plan.

On considère les points
\[
A(-1,2),\qquad B(5,0),\qquad C(3,4).
\]

\begin{enumerate}[label=\textbf{\arabic*.}]
    \item Calculer les longueurs \(AB\), \(AC\) et \(BC\). \barq{1.5}

    \item Le triangle \(ABC\) est-il rectangle ? isocèle ?  
    Toute réponse devra être démontrée. \barq{1.5}

    \item Déterminer une équation de la médiatrice de \([AB]\). \barq{1.5}

    \item Déterminer une équation de la droite \((BC)\). \barq{1}

    \item Montrer que la médiatrice de \([AB]\) et la droite \((BC)\) sont sécantes, puis déterminer leur point d'intersection \(I\). \barq{1.5}

    \item Interpréter géométriquement le fait que \(I\) appartienne à la médiatrice de \([AB]\). \barq{0.5}

    \item Soit \(M(x,y)\) un point du plan vérifiant
    \[
    \vv{MA}\cdot\vv{MB}=0.
    \]
    Montrer que cette condition équivaut à une équation cartésienne de cercle. \barq{2}

    \item Déterminer le centre et le rayon de ce cercle. \barq{1}

    \item Retrouver alors un résultat classique de géométrie sur l'angle \(\widehat{AMB}\). \barq{1}

    \item Déterminer les points \(M\) de la droite \((AC)\) vérifiant
    \[
    \vv{MA}\cdot\vv{MB}=0.
    \]
    \barq{2}

    \item On considère l'ensemble
    \[
    \Gamma=\left\{M(x,y)\in\R^2 \,\middle|\, MA^2-MB^2=10\right\}.
    \]
    Montrer que \(\Gamma\) est une droite. \barq{1.5}

    \item Donner une interprétation géométrique de ce lieu. \barq{0.5}

    \item On considère enfin l'ensemble
    \[
    \Sigma=\left\{M(x,y)\in\R^2 \,\middle|\, MA^2+MB^2=26\right\}.
    \]
    Montrer que \(\Sigma\) est soit vide, soit réduite à un point, soit un cercle.  
    Dans le cas présent, déterminer précisément \(\Sigma\). \barq{2}
\end{enumerate}

\vspace{1.2em}

\begin{center}
\begin{tabular}{|m{0.35\textwidth}|m{0.2\textwidth}|m{0.2\textwidth}|}
\hline
\centering \textbf{Partie} & \centering \textbf{Barème} & \centering \textbf{Note} \tabularnewline
\hline
\centering Problème 1 & \centering 24 & \tabularnewline
\hline
\centering Problème 2 & \centering 16 & \tabularnewline
\hline
\centering \textbf{Total} & \centering \textbf{40} & \tabularnewline
\hline
\end{tabular}
\end{center}

\vspace{1.2em}

\begin{center}
    \rule{0.6\textwidth}{0.5pt}\\[0.4em]
    \textit{Fin du devoir}\\[0.2em]
    \rule{0.6\textwidth}{0.5pt}
\end{center}

\end{document}